% =========================================================
% (2) GLOBAL COLOR DEFINITIONS
% =========================================================
% Semantic color system:
%
% Colors encode dynamical meaning rather than aesthetics:
% - Stability classes: Stable / Unstable / Saddle / Focus
% - Equilibrium families: HHEqui, HLEqui, MEqui, LLEqui, TEqui
%
% This allows model-independent visualization where
% color = semantic state, not arbitrary styling.
% =========================================================

% --- stability palette ---
\definecolor{Stable}{RGB}{0,110,0}
\definecolor{Unstable}{RGB}{255,77,77}
\definecolor{Saddle}{RGB}{255,165,0}
\definecolor{Focus}{RGB}{102,102,255}
\definecolor{Periodic}{RGB}{26,26,255}

% --- equilibrium-type palette ---
\definecolor{HHEqui}{RGB}{0,114,178}
\definecolor{HLEqui}{RGB}{230,159,0}
\definecolor{MEqui}{RGB}{0,158,115}
\definecolor{LLEqui}{RGB}{204,121,167}
\definecolor{TEqui}{RGB}{213,94,0}

% =========================================================
% (3) ATLAS TEMPLATE SYSTEM (GEOMETRIC DEFINITIONS)
% =========================================================
%
% Each \initAtlas defines a reusable geometric layout:
% - region identifiers (e.g., LL, HL, M, HH, C)
% - TikZ drawing instructions per region
% - semantic mapping to colors via \baColor{<key>}
%
% This separates:
%   geometry (structure)
%   from
%   data (region labeling)
% =========================================================

% ---------------------------------------------------------
% TEMPLATE: Cartel (5-region structured system)
% ---------------------------------------------------------
\initAtlas{Cartel}{LL,HL,M,HH,C}{
  % base blocks (semantic color injection)
  \fill[\baColor{LL}] (0,0) rectangle (1,1);
  \fill[\baColor{HL}] (1,0) rectangle (2,1);
  \fill[\baColor{M}]  (0,1) rectangle (1,2);
  \fill[\baColor{HH}] (1,1) rectangle (2,2);

  % refinement region (high-resolution structure)
  \fill[\baColor{C}] (0,0) rectangle (0.5,0.5);

  % structural grid overlay
  \draw[line width=2pt,black] (0,2)--(0,0)--(2,0);
  \draw (0,1)--(2,1);
  \draw (1,0)--(1,2);

  % refinement grid
  \draw (0,0.5)--(0.5,0.5);
  \draw (0.5,0)--(0.5,0.5);

  % highlight boundary
  \draw[line width=2pt,red] (0,0.04)--(0,2);
}

% ---------------------------------------------------------
% TEMPLATE: SimpleState (2-region system)
% ---------------------------------------------------------
\initAtlas{SimpleState}{SL,SR}{
  \fill[\baColor{SL}] (0,0) rectangle (1,2);
  \fill[\baColor{SR}] (1,0) rectangle (2,2);
}

% ---------------------------------------------------------
% TEMPLATE: Circles
% ---------------------------------------------------------
\initAtlas{Circles}{NorthWest,NorthEast,SouthWest,SouthEast,Center}{
  \fill[\baColor{NorthWest}] (-1, 1) circle (0.35);
  \fill[\baColor{NorthEast}] ( 1, 1) circle (0.35);
  \fill[\baColor{SouthWest}] (-1,-1) circle (0.35);
  \fill[\baColor{SouthEast}] ( 1,-1) circle (0.35);
  \fill[\baColor{Center}]    ( 0, 0) circle (0.45);

  \node at (-1, 1) {\tiny NW};
  \node at ( 1, 1) {\tiny NE};
  \node at (-1,-1) {\tiny SW};
  \node at ( 1,-1) {\tiny SE};
  \node at ( 0, 0) {\tiny C};

  \draw (-1, 1) circle (0.35);
  \draw ( 1, 1) circle (0.35);
  \draw (-1,-1) circle (0.35);
  \draw ( 1,-1) circle (0.35);
  \draw ( 0, 0) circle (0.45);
}

